\subsubsection{Mission timelines and the structure of the energy budget}
Figure~\ref{fig:X9} shows the power drawn by each aircraft over time for one wind realisation. The aerodynamic benefit is invisible at this scale: the followers' transit power differs from the leader's by $\approx$\SI{8}{W} on \SI{143}{W}, and the dominant structures are the mapping legs and the turns (the Dubins loops of the 17-strip plan are visible as the orange blocks between mapping legs). The strategies with formation differ from S1 by the reconfiguration blocks (violet) at the start, the split, the merge and the break-up, and by the waiting blocks (magenta) of the aircraft that finish early. This is the structural reason that the formation overheads, small in isolation, add up to more than the aerodynamic saving at \SI{2.5}{km}.
\begin{figure}[H]\centering\includegraphics[width=\textwidth]{X9_power_timeline}
\caption{Load power versus time for the three aircraft (rows) under S1, S3 and S4 (columns) for one environment realisation (wind \SI{3.5}{m/s}); colours are mission phases.}\label{fig:X9}\end{figure}

\subsubsection{Leader rotation and battery balance}
Because only the followers profit from the wake, a fixed leader would carry the full cost of the formation phases. Rotating the role cyclically makes the formation-phase energy identical for the three aircraft; without rotation the leader would spend about \SI{4}{Wh} more than a follower over the S3 transit (Fig.~\ref{fig:X10}b), which is \SI{10}{\percent} of the aircraft's mission energy. The maximum per-aircraft energy of S3 and S4 is \nmisEmaxSThree\ and \nmisEmaxSFour\,Wh, against \nmisEmaxSOne\,Wh for S1; the usable battery (\SI{177.6}{Wh}) is never approached by any strategy except the rigid V, whose per-aircraft energy is \nmisEmaxSTwo\,Wh.
\begin{figure}[H]\centering\includegraphics[width=\textwidth]{X10_leader_rotation_energy}
\caption{(a) Per-aircraft energy (maximum, filled circle, and minimum, open circle, over the three aircraft; mean $\pm$ SD over trials) for four strategies, against the usable battery energy; (b) leader-minus-follower energy gap over the S3 formation transit that cyclic rotation eliminates.}\label{fig:X10}\end{figure}

\subsubsection{Monte-Carlo diagnostics}
The paired difference S3$-$S1 has a mean of \nmcDiffMean\,Wh and a trial-to-trial SD of only \nmcDiffSD\,Wh (Fig.~\ref{fig:X8}a): the running mean has converged to within a few hundredths of a watt-hour after $\approx$40 trials, and the bootstrap interval is narrower than the line width beyond 80. By contrast, the energy of a single strategy varies by \nmcRange\,Wh over the wind range (SD \nmcSDOne\,Wh); a linear regression on wind speed alone explains $R^2=\nmcRTwo$ of that variance with a slope of \nmcSlope\,Wh per m/s (Fig.~\ref{fig:X8}b), i.e.\ the strip headwind/tailwind asymmetry and the longer ground-speed-limited legs. The strategy \emph{difference} depends only weakly on the wind: a harmonic regression on wind direction explains $R^2=\nmcDirRTwo$ of its variance (Fig.~\ref{fig:X8}c). Tight paired intervals here therefore document simulation determinism, not field repeatability.
\begin{figure}[H]\centering\includegraphics[width=\textwidth]{X8_monte_carlo_diagnostics}
\caption{(a) Running mean of the S3$-$S1 difference with bootstrap 95\,\% interval versus number of trials; (b) S1 energy versus wind speed (colour: wind direction); (c) S3$-$S1 versus wind direction (polar).}\label{fig:X8}\end{figure}
